\documentclass[11pt]{article}
\usepackage[utf8]{inputenc}
\usepackage[T1]{fontenc}
\usepackage{lmodern}
\usepackage[margin=1in]{geometry}
\usepackage{amsmath,amssymb}
\usepackage{booktabs}
\usepackage{longtable}
\usepackage{hyperref}
\usepackage{url}
\usepackage{enumitem}
\usepackage{xcolor}
\hypersetup{colorlinks=true, linkcolor=blue!50!black, citecolor=blue!50!black,
            urlcolor=blue!50!black}

\title{\textbf{The Agentic Internet Needs a Trust Layer}\\[0.3em]
\large XCP: Verifiable Delegation, Graded Reachability, and
Evidence-Conditioned Settlement for the Agentic Economy}

\author{Anonymous Contributors\\
\small AGI-Gateway / XCP Project\\
\small \texttt{https://github.com/AGI-Gateway/XCP}}

\date{Working paper --- draft for comment}

\begin{document}
\maketitle

\begin{abstract}
The agentic web is technically feasible and economically imminent, but it is
being built on an authorization model designed for humans clicking buttons.
Agents today act with borrowed credentials, unbounded scope, and no portable
evidence that anything was authorized. We argue that the missing layer is not
another transport protocol but a \emph{trust layer}: a way for one party to
verify, per request, \emph{who} is acting, \emph{on whose authority},
\emph{within what bounds}, and \emph{with what recourse}.

We present XCP, a stateless trust layer for agentic infrastructure. XCP
contributes five mechanisms: (i) a \emph{trust lattice} that composes an agent's
accountability with its human principal's verification, and which narrows but
never widens authority; (ii) \emph{per-request credentials} bound to one method,
one tool, one argument digest, and a seconds-long expiry, adapted to MCP
2026-07-28's removal of protocol sessions; (iii) a \emph{trust firewall} that
grades reachability rather than gating it binary, so an unknown server is
observable but never binding; (iv) \emph{evidence-conditioned settlement}, where
payment releases against a verifiable receipt rather than an assertion of
delivery; and (v) \emph{federation on stable identity}, so that certificate
rotation does not orphan peers.

We report a working reference implementation: 485 tests, a black-box conformance
suite that certifies third-party implementations, a measured trust-layer overhead
of \textbf{0.21\,ms p50} over a bare proxy --- roughly $0.7\%$ of a single
$30$\,ms cross-region round trip --- and a catalog of 9{,}576 endpoints whose
trust classes are, deliberately, mostly \texttt{unknown}. We are explicit about
what is unproven: no independent security audit, an unaudited and undeployed
contract, and a federation that has been demonstrated between two nodes rather
than at scale.

We close by arguing that the technical layer is necessary but not sufficient.
Drawing on recent work in normative and legal infrastructure for the agentic web,
we sketch a research agenda for \emph{Agentic Policy} and \emph{Agentic Society}:
how verifiable delegation might become a substrate for augmenting human
collective intelligence rather than displacing it, and what ethical first
principles should govern an autonomous future in which life and consciousness are
meant to thrive.
\end{abstract}

\section{Introduction}

Two claims are now widely accepted and jointly uncomfortable. The first is that
AI agents acting on a principal's behalf across the open web are technically
feasible today~\cite{pattison2026}. The second is that the infrastructure
governing what those agents may do was designed for a different
world~\cite{yang2025,zhu2026ioai}.

The gap is not connectivity. The Model Context Protocol (MCP), Agent-to-Agent
(A2A) protocols, and adjacent specifications have largely solved \emph{how} an
agent reaches a tool. What remains unsolved is \emph{whether it should}, and how
a counterparty could tell. Zhu's layered analysis of the Internet of Agentic
AI names this precisely: connectivity, identity and trust, capability discovery,
task semantics, governance, and an economic layer are distinct requirements, and
``no single protocol can solve interoperability by itself''~\cite{zhu2026ioai}.
Today's deployments have the first and, increasingly, the third. The middle
layers are improvised.

The improvisation has a shape. An agent is handed a long-lived API key scoped to
everything that key can do. It presents no evidence of which human authorized it,
under what limit, for how long. When it acts, no artifact survives that a third
party could check. When something goes wrong, the only remedy is to revoke the
key --- which revokes every agent using it. This is not a protocol gap; it is a
\emph{trust} gap, and it is the reason platforms increasingly treat all automated
traffic as hostile, blocking legitimate delegated agents alongside malicious
bots~\cite{pattison2026}.

\paragraph{Contribution.} This paper presents XCP, a trust layer designed to sit
between an agent and the tools it calls. Our contributions are:

\begin{enumerate}[leftmargin=1.4em,itemsep=0.15em]
\item A \textbf{trust lattice} (\S\ref{sec:lattice}) composing agent
      accountability with human verification, with a tested invariant that
      enterprise entitlements may only \emph{narrow} authority, never widen it.
\item \textbf{Stateless per-request credentials} (\S\ref{sec:stateless}) adapted
      to MCP 2026-07-28, bound to method, tool, argument digest, TLS footprint,
      and a seconds-long expiry --- so a stolen credential is worth
      substantially less than a stolen session.
\item A \textbf{trust firewall} (\S\ref{sec:firewall}) that grades reachability
      across four server classes, producing obligations (sandbox, quarantine,
      receipt) that travel with the decision.
\item \textbf{Evidence-conditioned settlement} (\S\ref{sec:receipts}): an
      append-only call chain and signed receipt that make payment conditional on
      verifiable delivery rather than asserted delivery.
\item \textbf{Federation on stable identity} (\S\ref{sec:federation}), where a
      node's identity is a long-lived key and the certificate is bound
      underneath it, so renewal does not orphan peers or invalidate
      attestations.
\end{enumerate}

We also report measurements (\S\ref{sec:eval}), state limitations plainly
(\S\ref{sec:limits}), and propose a normative research agenda
(\S\ref{sec:future}).

\section{Background and Related Work}

\paragraph{The agentic web as an economic substrate.}
Yang et al.~\cite{yang2025} characterize the Agentic Web across three
dimensions --- intelligence, interaction, and economics --- and identify agent
discovery, trust, billing, and governance as open infrastructure problems. Their
framing motivates ours: if agents are to transact, the economic dimension cannot
be bolted onto a substrate that cannot say who acted. XCP addresses the trust and
accounting dimensions specifically, and treats discovery as a solved-enough
problem that should not be re-solved on-chain.

\paragraph{Internet-scale agent coordination.}
Zhu's Internet of Agentic AI (IoAI)~\cite{zhu2026ioai} provides the most direct
architectural antecedent. It enumerates the threats XCP is built against ---
impersonation, Sybil attacks, registry poisoning, unauthorized delegation,
cascading failures --- and argues that existing human-centric identity systems
are insufficient for ephemeral, autonomous, cross-organizational agents. Its
layered interoperability table maps closely onto XCP's design: we implement the
identity-and-trust, governance, and economic rows, and deliberately do not
reimplement connectivity or capability discovery. Zhu's companion work on the
Internet of Agentic Things (IoAT)~\cite{zhu2026ioat} extends the same argument to
cyber-physical settings, where a reasoning failure becomes a physical action --
raising the stakes on exactly the authorization and reversibility distinctions
XCP encodes.

\paragraph{Normative infrastructure.}
Pattison, Boulos, Kolt, Li, Piccardi, and Lazar~\cite{pattison2026} argue that
the bottleneck to the agentic web is no longer technical but normative, and
propose a triad of \emph{delegation} (a user entitled to access should be able to
exercise it through an authorized agent), \emph{transparency} (agents identify
themselves and their principals; platforms disclose how they treat agents), and
\emph{proportional restriction} (restrictions address concrete harms by the least
restrictive means).

This triad is, to our reading, the clearest statement of what a trust layer
should make \emph{checkable}. XCP's mandate gate is a delegation mechanism;
\texttt{XCP-Agent-Identity} and the verified-by header implement bidirectional
transparency; and graded reachability is a proportionality mechanism --- an
unknown server is observed rather than blocked, and blocked only where a concrete
risk is established. We regard~\cite{pattison2026} as the normative specification
for which XCP is one candidate implementation, and we note their observation that
agents should be able to identify themselves \emph{as authorized} without
necessarily identifying the individual principal --- a property our lattice
supports via tier disclosure without identity disclosure.

\paragraph{Legal and adjudicative infrastructure.}
Chaffer et al.~\cite{chaffer2026} argue for distributed legal infrastructure for
a trustworthy agentic web, spanning identity, machine-readable constraints,
adjudication, market policy, and portability. Their framing clarifies the
boundary of our contribution: XCP produces the \emph{evidence} that an
adjudicative layer would need --- tamper-evident call chains, signed receipts,
per-request credentials --- but does not itself adjudicate. We regard evidence
production and dispute resolution as separable, and deliberately implement only
the former.

\paragraph{Positioning.}
Relative to this literature, XCP is narrower and more concrete. It does not
propose a governance regime~\cite{pattison2026,chaffer2026}, a full coordination
architecture~\cite{zhu2026ioai}, or an economic model~\cite{yang2025}. It
proposes a mechanism layer that those agendas presuppose, and reports what
happens when you actually build it.

\section{The Trust Gap}

We characterize the gap concretely, since the fix follows from the diagnosis.

\paragraph{Credentials carry no principal.}
An API key says \emph{which account} is acting, never \emph{which human
authorized this call}. At the moment an agent spends money or touches regulated
data, the accountable party is unrepresented in the request.

\paragraph{Scope is ambient, not per-call.}
OAuth scopes are granted at consent time and apply to every subsequent call.
Nothing binds \emph{this} call to \emph{that} authorization. An agent that was
authorized to read a calendar and later writes to it presents identical
credentials.

\paragraph{Sessions are worth stealing.}
Session identifiers are bearer tokens with broad scope and long life. MCP
2026-07-28's removal of \texttt{Mcp-Session-Id} is, in our reading, an
opportunity rather than a loss: it forces each request to be self-describing,
which is the precondition for making a credential cheap to steal.

\paragraph{Delivery is asserted, not proven.}
When an agent pays for work, it has the provider's word that work occurred. There
is no artifact a third party can verify offline, which makes automated settlement
between strangers a matter of trust rather than evidence.

\paragraph{Reachability is binary.}
Access control answers allow/deny. But an agent encountering an unknown MCP
server does not need a binary answer; it needs to know that it may \emph{read}
the server's output without treating it as authoritative, may not let it
\emph{write}, and must quarantine what it returns. Binary gating forces a choice
between paralysis and recklessness.

\section{Design}

XCP is organized as three planes --- identity, governance, payments --- over a
stateless request model. We describe each mechanism and the invariant it
enforces.

\subsection{The Trust Lattice}\label{sec:lattice}

Authority is composed from two independent axes rather than assigned as a role.

\begin{center}
\small
\begin{tabular}{@{}llll@{}}
\toprule
& \textbf{H0} public & \textbf{H1} general SSO & \textbf{H2} enterprise IdP \\
\midrule
\textbf{A0} free/ephemeral   & sandbox      & sandbox     & shadowed \\
\textbf{A1} registry-backed  & observe only & consumer commerce & constrained \\
\textbf{A2} company-backed   & principal    & delegated   & full settlement \\
\bottomrule
\end{tabular}
\end{center}

The ceiling is the weaker leg, with one deliberate exception: A2$\times$H0, where
the company is itself the accountable principal and no human verification is
required. The lattice is not decorative --- it drives mandate templates, rail
eligibility, receipt requirements, and (as we discuss in \S\ref{sec:abuse})
resource quotas.

\paragraph{Invariant: entitlements narrow only.}
When an enterprise IdP asserts group membership, cost centre, or approval limit,
those claims may \emph{reduce} what a mandate permits and may never expand it. A
compromised or overreaching IdP therefore cannot manufacture authority. This is
enforced in code and tested, because it is the property on which enterprise
adoption depends: identity is not authority.

\subsection{Stateless Per-Request Credentials}\label{sec:stateless}

MCP 2026-07-28 removed protocol sessions and made \texttt{Mcp-Method} and
\texttt{Mcp-Name} mandatory headers. XCP treats this as the authorization
primitive. A request credential is an EIP-712 signature binding:

\[
\texttt{agent} \;\|\; \texttt{chain} \;\|\; \texttt{tier} \;\|\;
\texttt{method} \;\|\; \texttt{tool} \;\|\; H(\texttt{args}) \;\|\;
\texttt{footprint} \;\|\; \texttt{exp}
\]

with a default expiry of 60 seconds. Verification is stateless: a cold replica
can validate without shared state. Scope derives from headers alone
(\texttt{scope\_for(method, name)}), so a gateway can make an authorization
decision at header speed --- measured at $0.27\,\mu$s --- before parsing a body.

The security property is comparative rather than absolute. A stolen credential
buys one tool call, on one argument set, for under a minute, over the same TLS
channel. Replaying it against a different tool fails on scope mismatch; we test
this directly.

\subsection{Trust Firewall: Graded Reachability}\label{sec:firewall}

Servers are classified into four grades --- \textsc{unknown}, \textsc{probed},
\textsc{attested}, \textsc{contracted} --- and the decision is a function of
caller tier $\times$ server class, producing an effect
(\textsc{deny}/\textsc{observe}/\textsc{allow}) plus \emph{obligations} that
travel with it.

Four rules are enforced and tested:

\begin{itemize}[leftmargin=1.4em,itemsep=0.1em]
\item An \textsc{unknown} server is always \textsc{observe}: read-only,
      sandboxed, output quarantined, nothing binding.
\item A free agent (A0) may read the world but may not write to it, at any
      server class.
\item Settlement requires \textsc{attested} or better \emph{and} an accountable
      caller.
\item Output below \textsc{contracted} is always quarantined.
\end{itemize}

This is our reading of proportional restriction~\cite{pattison2026} in
mechanism: the default is not exclusion but bounded engagement, and restriction
escalates with demonstrated risk rather than with unfamiliarity.

\subsection{Evidence-Conditioned Settlement}\label{sec:receipts}

A \texttt{TaskSpec} commits to the work before it begins. A \texttt{CallChain}
accumulates an append-only hash chain, $h_i = H(h_{i-1} \| d_i)$, over the calls
actually made. A signed \texttt{Receipt} binds the two into an evidence bundle
verifiable offline by a party who was not present.

Escrow then moves \textsc{open} $\rightarrow$ \textsc{delivered} $\rightarrow$
\textsc{accepted} $\rightarrow$ \textsc{released}. Auto-accept is off by default,
and \emph{delivery alone does not release funds}. Every settling tier requires a
receipt --- a tested invariant, because the alternative is an economy in which
payment is conditioned on assertion.

\subsection{Federation on Stable Identity}\label{sec:federation}

An early version of this design made node identity the certificate footprint,
$\mathrm{keccak256}(\mathrm{DER}(\mathrm{cert}))$. This was a mistake worth
reporting: it makes identity an \emph{ephemeral credential}. With 60--90 day
certificate renewal, every peer sees a stranger four times a year, attestations
others issued become worthless, and revocations key to an identity that no longer
exists. A trust network whose members lose their identity quarterly is not a
trust network.

The correction separates the two: a node's identity is a long-lived key, and the
certificate is bound underneath it by a signed statement carrying a monotonic
sequence number and validity window. Renewal publishes a new \emph{binding}, not
a new node. Two properties are required to make this safe and are tested:
\emph{overlap} (the previous certificate is accepted within a grace window, since
both are briefly live during rotation) and \emph{anti-downgrade} (a peer refuses
any binding whose sequence is not greater than the newest it has seen, or an
attacker could replay a superseded binding to reinstate a rotated-away key).

Federation is otherwise deliberately minimal: peer verification is three local
checks with no third party, trust decays $0.5$ per hop capped at three hops,
revocation gossips only from already-verified peers, and nodes may disagree
without consensus.

\subsection{Abuse Controls: Capacity Follows Authority}\label{sec:abuse}

A node published in a discovery catalog, routing for callers it has never met,
with no rate limit, is an open relay. We treat the lattice as a capacity
allocation as well as an authority allocation: an anonymous caller receives 20
cost units per minute and one concurrent request; an A2$\times$H2 caller receives
30{,}000 and 64.

Four limits are enforced because the failure modes differ --- rate (flooding),
burst (legitimate spikes), concurrency (slow-loris, which a requests/second limit
does not touch), and body size --- plus a global ceiling that binds regardless of
tier, since a node that dies serving trusted traffic is equally down. Charges are
cost-weighted: a multi-hop chain costs twelve units where a \texttt{tools/list}
costs one.

Two implementation notes generalize. First, the limiter fails \emph{closed}:
failing open converts any limiter bug into an abuse bypass. Second, the
tracked-caller table is a bounded LRU, because a limiter keyed on an
attacker-supplied identifier is itself a memory-exhaustion vector.

\subsection{Credential Custody and Erasure}

Two mechanisms address the fact that a trust layer is also a data-handling layer.

\paragraph{Sealed credentials.} A hosted API wrapper needs the caller's upstream
credential. Forwarding it through the gateway means the gateway operator reads
every user's API key --- disqualifying for a trust layer. Instead the caller
seals the credential to the \emph{wrapper's} public key (ephemeral X25519,
HKDF-SHA256, AES-256-GCM, with the recipient key and operation context in the
AAD). The gateway forwards ciphertext it has no key for. The guarantee is narrow
and we state it precisely: this removes the \emph{gateway} operator from the
trust set, not the wrapper operator.

\paragraph{Crypto-shredded erasure.} The call chain is designed so that removing
a record changes the root --- which is exactly what makes it undeletable, and
collides with the right to erasure. Deleting records destroys the integrity
property; refusing erasure because a structure is append-only is not a lawful
basis. We resolve this by never placing personal data \emph{in} the chain:
payloads are encrypted under a per-subject key and the chain commits to the
ciphertext digest. Erasure destroys the key. Content becomes unreadable ---
including to the operator --- while the chain still verifies and its root is
unchanged. We note that a hash of personal data remains personal data when
linkable, so digesting a field does not discharge an erasure obligation; key
destruction does.

\section{Implementation and Evaluation}\label{sec:eval}

The reference implementation is open source (MIT, including the specification)
and comprises a gateway, reference server, verifier, clients, a federation node,
and a CLI. It targets MCP 2026-07-28 and interoperates with ERC-8004 identity
registries where present, while functioning with no chain dependency at all.

\paragraph{Conformance.}
Because a federation is a protocol rather than a program, we ship a black-box
conformance suite: it speaks HTTP, imports nothing from the implementation under
test, and can certify a Go, Rust, or TypeScript gateway identically. Eighteen
clauses span core, stateless, security, and federation profiles, with MUST/SHOULD
separation. Most checks are \emph{negative} --- a suite that only exercised happy
paths would certify nothing, since a gateway that accepts everything passes every
happy-path test. We validate the discrimination: the reference implementation
certifies 15/15, while a deliberately permissive gateway fails seven clauses.

\paragraph{Performance.}
We measure against a baseline passthrough proxy with an identical upstream, since
an absolute latency figure is not a claim about anything.

\begin{center}
\small
\begin{tabular}{@{}lrrr@{}}
\toprule
& p50 & p95 & p99 \\
\midrule
Baseline proxy (no trust layer)        & 1.91\,ms & 2.11\,ms & 2.36\,ms \\
XCP gateway (enforce, all controls)    & 2.12\,ms & 2.38\,ms & 2.55\,ms \\
\textbf{Trust-layer overhead}          & \textbf{+0.21\,ms} & \textbf{+0.27\,ms} & \\
\bottomrule
\end{tabular}
\end{center}

The overhead is $1.11\times$ baseline, or roughly $0.7\%$ of a single $30$\,ms
cross-region round trip. A refused call costs $0.59$\,ms and is metered
pre-authentication, so an unauthenticated flood remains bounded. Signature
verification dominates at $185\,\mu$s; scope derivation ($0.27\,\mu$s), lattice
resolution ($6\,\mu$s), and rate limiting ($2.6\,\mu$s) are negligible by
comparison.

Two measurement findings are worth reporting because both are easy to
misattribute. First, an earlier measurement showed $24$\,ms of overhead; the
cause was the absence of upstream connection pooling, not verification work.
$99\%$ of observed latency had nothing to do with the trust layer. Second,
\texttt{eth-keys} names its \emph{pure-Python} ECDSA implementation
\texttt{NativeECCBackend}; without native \texttt{libsecp256k1} bindings,
verification costs $6{,}576\,\mu$s instead of $185\,\mu$s --- a $35\times$ penalty
a deployment pays silently. We now pin the dependency, warn at startup, and
expose the active backend in \texttt{/health}.

\paragraph{Catalog.}
We ship a catalog of 9{,}576 endpoints across three kinds: 158 \emph{routable}
(remote MCP endpoints reachable today), 2{,}528 \emph{wrappable} (public REST
APIs with OpenAPI specifications, one generator step from routable), and 6{,}869
\emph{installable} (packages with nothing to route to until an operator runs
one). The distinction is enforced rather than cosmetic: marking a REST API
routable would send agents to endpoints where every call fails.

Critically, verification caps trust: the shipped catalog never claims
\textsc{attested} or \textsc{contracted}, and of the 158 routable entries only
five are vendor-confirmed. A catalog that asserted its own trustworthiness would
be the attack.

\section{Limitations}\label{sec:limits}

We state these plainly because a security claim that overstates itself is a
defect.

\begin{itemize}[leftmargin=1.4em,itemsep=0.15em]
\item \textbf{No independent security audit.} The implementation is tested (485
      tests across 21 suites) and enforcement contains no stubs, but no external
      party has reviewed it.
\item \textbf{The contract is unaudited and undeployed.} Compiling and executing
      \texttt{NodeRegistry.sol} for the first time surfaced three real defects
      --- a griefing vector allowing a stranger to poison a route identifier, an
      unaccounted forfeited-stake path, and the absence of any unbond mechanism,
      which made the bond economically unusable. All are fixed and tested. No
      formal verification, fuzzing, or economic modelling has been performed, and
      the reward parameters are guesses.
\item \textbf{Federation is demonstrated, not proven at scale.} Two nodes have
      been run against each other, which surfaced four defects --- including that
      the federation module had no reference in the gateway at all, existing as a
      tested library no running node could reach. Behaviour at hundreds of nodes
      is unknown.
\item \textbf{Optimistic claiming assumes someone challenges.} The fraud proofs
      are sound but only bite if a party with a stake actually checks. In a
      sparse federation that assumption is weak.
\item \textbf{Measurements are loopback and Python.} The relative overhead is
      flattered by the absence of network latency and penalized by the reference
      language; a Go or Rust implementation would move the CPU-bound numbers.
\item \textbf{Erasure does not reach peers or backups.} Key destruction removes a
      key from a process. Copies already federated are beyond reach.
\item \textbf{Compliance mapping is not compliance.} We map controls to the EU AI
      Act, DORA, GDPR, NIS2, SOC 2 and others with an explicit gap register, and
      no lawyer has reviewed it.
\end{itemize}

\section{Future Work: Agentic Policy and Agentic Society}\label{sec:future}

The mechanisms above answer a narrow question --- \emph{can a counterparty verify
this action was authorized?} --- and we think that question is prior to the
interesting ones. We sketch four directions, in increasing order of difficulty
and decreasing order of our confidence.

\subsection{From Verifiable Delegation to Agentic Policy}

Pattison et al.~\cite{pattison2026} observe that platforms and agent developers
are locked in an arms race that burns societal value, largely because no shared
vocabulary distinguishes a user-authorized agent from an indiscriminate scraper.
A trust layer supplies that vocabulary, but does not by itself create the
incentive to use it.

The open problem is \emph{policy expressed as mechanism}. If proportional
restriction is to be more than an aspiration, a platform needs to state its
access policy in a form an agent can read, satisfy, and be held to --- and a
regulator needs to be able to check that the stated policy is the enforced one.
XCP's graded reachability and mandate proofs are candidate primitives; what is
missing is a machine-readable policy language with adjudicative standing, which
is where the legal-infrastructure agenda~\cite{chaffer2026} and the mechanism
agenda must meet.

\subsection{Economic Fairness Under Agentic Inequality}

Delegated agents are not equally capable, and early evidence suggests
better-resourced agents extract predictable gains from weaker
counterparties~\cite{pattison2026}. Evidence-conditioned settlement makes
transactions \emph{verifiable}; it does nothing to make them \emph{fair}.

We regard this as the most urgent economic question the agentic economy poses.
If agent-mediated commerce simply re-encodes existing resource asymmetries at
machine speed, the promise that agents will democratize access to expertise
inverts into a mechanism for concentrating advantage. Candidate directions
include capability-matched negotiation (agents transact only with counterparties
of comparable capability), baseline agent provision as a public good, and
settlement rules that make exploitative asymmetry detectable in the evidence
bundle itself. None of these are solved, and we do not claim XCP addresses them.

\subsection{Agentic Society: Governing Emergence}

Zhu frames the central design tension as \emph{controlled emergence}: harnessing
system-level behaviour arising from local interaction while preserving
predictability, alignment, and accountability~\cite{zhu2026ioai}. As agent
populations grow, the object of governance shifts from individual actions to
\emph{interaction patterns} --- collusion, cascading failure, coordinated
behaviour that no participant intended and none can unilaterally stop.

XCP's contribution here is negative but real: tamper-evident call chains make
interaction patterns \emph{observable after the fact}, and graded reachability
bounds how far an unknown participant's influence propagates. What is missing is
the positive capability --- detecting emergent misalignment while it is
happening, at population scale, without building the surveillance apparatus that
such detection naively implies. We note the tension directly: the same evidence
that enables accountability enables monitoring, and our privacy design (scrubbed
telemetry, crypto-shredded erasure, commitments that publish totals but never
counterparties) is an attempt to keep that boundary explicit rather than to
pretend it does not exist.

In cyber-physical settings this ceases to be abstract. As Zhu argues for the
Internet of Agentic Things, a reasoning failure becomes a physical
action~\cite{zhu2026ioat}, and an agent's authority to \emph{act} must therefore
be narrowed by reversibility class --- a setpoint change, a door unlock, and an
emergency shutdown should not share approval logic. XCP's obligations mechanism
is a plausible carrier for such distinctions; we have not built them.

\subsection{Ethical First Principles for an Autonomous Future}

We close with the commitment that motivated this work, stated as a position
rather than a result.

The purpose of building infrastructure for autonomous agents should be to
\emph{augment human collective intelligence}, not to substitute for it. This is a
design constraint with teeth, and it argues for specific properties:

\begin{description}[leftmargin=1.4em,itemsep=0.2em]
\item[Principals remain human.] Authority originates with a person or an
      accountable organization and is delegated, never generated. Our lattice
      invariant --- entitlements narrow, never widen --- is a small technical
      expression of this. Authority that an agent can manufacture is authority
      no human granted.
\item[Delegation is revocable and bounded.] Autonomy that cannot be withdrawn is
      not delegation but abdication. Short-lived credentials, per-call scope
      binding, and unsuppressable revocation are not merely security features;
      they are what keeps a human in the loop of consequence even when they are
      not in the loop of execution.
\item[Evidence belongs to the principal.] An audit trail that only the platform
      can read reproduces the asymmetry agents were supposed to dissolve.
      Receipts verifiable \emph{offline}, by a party who was not present, are the
      mechanism by which accountability flows toward the person rather than the
      intermediary.
\item[Pluralism over consensus.] We deliberately built a federation in which
      nodes may disagree and no global agreement is required. A trust
      infrastructure that forces one answer about who is trustworthy becomes a
      chokepoint, and chokepoints get captured.
\end{description}

Whether an autonomous future is one in which life and consciousness thrive
depends less on the capability of the agents than on whether the infrastructure
they run on preserves the standing of the beings they act for. We think that is
substantially an infrastructure question, which is why we have approached it as
one. But we hold the position loosely: mechanism can make good outcomes
\emph{possible} and bad outcomes \emph{detectable}; it cannot make good outcomes
\emph{happen}. That requires the societal conversation
that~\cite{pattison2026} calls for, and to which this paper is intended as a
contribution of one part --- the part where you can check the claims by running
the code.

\section{Conclusion}

The agentic economy is being built on an authorization model designed for humans
clicking buttons. We have argued that the missing layer is a trust layer, and
presented one: a lattice that composes agent accountability with human
verification, per-request credentials cheap enough to be worth little when
stolen, reachability that is graded rather than binary, settlement conditioned on
verifiable evidence, and federation that survives certificate rotation.

The measured cost is $0.21$\,ms. We think that is the wrong number to focus on.
The right one is that a counterparty can now answer, per request, who is acting,
on whose authority, within what bounds, and with what recourse --- and can check
the answer without trusting us.

\paragraph{Availability.}
Source, specification, conformance suite, benchmarks, and the full gap register
are at \url{https://github.com/AGI-Gateway/XCP} under MIT. We particularly invite
adversarial review of the trust lattice invariants, the contract, and the
conformance suite's negative checks --- and would rather receive a clause that
the reference implementation \emph{fails} than one it passes.

\paragraph{Status.}
XCP and ERC-8004x are draft proposals originating in this project. MCP, A2A,
ERC-8004, and the frameworks cited here are independent work by others.

\begin{thebibliography}{9}

\bibitem{yang2025}
Y. Yang et al.
\newblock Agentic Web: Weaving the Next Web with AI Agents.
\newblock \emph{arXiv:2507.21206}, 2025.

\bibitem{chaffer2026}
T. J. Chaffer et al.
\newblock Distributed Legal Infrastructure for a Trustworthy Agentic Web.
\newblock \emph{arXiv:2603.06884}, 2026.

\bibitem{pattison2026}
C. Pattison, M. Boulos, N. Kolt, C. Li, T. Piccardi, and S. Lazar.
\newblock The Agentic Web Requires New Normative Infrastructure.
\newblock \emph{arXiv:2606.10711}, 2026.

\bibitem{zhu2026ioai}
Q. Zhu.
\newblock The Internet of Agentic AI: Communication, Coordination, and
Collective Intelligence at Scale.
\newblock \emph{arXiv:2606.12835}, 2026.

\bibitem{zhu2026ioat}
Q. Zhu.
\newblock Internet of Agentic Things: Networked AI Agents for Closed-Loop IoT
Orchestration.
\newblock \emph{arXiv:2607.12662}, 2026.

\end{thebibliography}

\end{document}
